\section{Mathematical Preliminaries and Problem Formulation}
\label{problem}
To conceptualize the virtual encoder system, the kinematic inaccuracies of a Cartesian XZ-head/Y-bed architecture must be mathematically modeled. Machine tool calibration requires the extraction and compensation of three primary error components: linear step drift, mechanical backlash, and structural axis skewness. 

\subsection{Linear Step Derivation}
The theoretical linear motion resolution of a pulley-and-belt mechanism is dictated by the motor's step angle, microstepping configurations, and the physical parameters of the transmission components \cite{M92}. The nominal configuration governs the steps executed per millimeter ($Steps/mm_{default}$) via the following analytical relation:

\begin{equation}
Steps/mm = \frac{\theta_{rev} \times M_f}{P_b \times N_g}
\end{equation}

where $\theta_{rev}$ is the physical steps per full motor revolution (typically 200 steps for a $1.8^\circ$ stepper motor), $M_f$ is the electronic microstepping factor applied by the driver, $P_b$ represents the belt pitch (e.g., 2~mm for standard GT2 profiles), and $N_g$ denotes the tooth count of the drive pulley. Mechanical degradation and wear alter this theoretical framework over time, creating a discrepancy between commanded distances ($D_{cmd}$) and physical displacements ($D_{act}$). The rectified steps per millimeter profile ($Steps/mm_{new}$) is computed via a proportional scale matching law:

\begin{equation}
Steps/mm_{new} = \left(\frac{D_{cmd}}{D_{act}}\right) \times Steps/mm_{old}
\end{equation}

\subsection{Backlash Modeling}
Mechanical backlash manifests as a nonlinear deadzone error when an axis undergoes a reversal of motion direction. The lost motion arises from physical clearances between meshing transmission components (e.g., belt-pulley slack) \cite{papageorgiou2019}. Mathematically, if $x_{in}(t)$ represents the commanded actuator position and $x_{out}(t)$ defines the actual joint coordinate, the structural deadzone is formulated as:

\begin{equation}
x_{out}(t) = 
\begin{cases} 
x_{in}(t) - d_b & \text{if } \dot{x}_{in}(t) > 0 \\
x_{in}(t) + d_b & \text{if } \dot{x}_{in}(t) < 0 
\end{cases}
\end{equation}

where $d_b$ represents the absolute backlash width. The virtual encoder quantifies this phenomenon by monitoring positional states across a predefined bi-directional travel sequence, evaluating the absolute divergence between the original anchor coordinate and the coordinates achieved post-reversal.

\subsection{Axis Skewness Formulation}
Structural misalignment during assembly or asymmetric belt tensioning destroys the strict orthogonality between the X and Y movement vectors ($\vec{V}_x$ and $\vec{V}_y$). The non-orthogonality angle ($\theta_{skew}$) represents the deviation of the internal angle ($\theta_{int}$) from an ideal $90^\circ$ perpendicular intersection. Let the absolute orientation angles ($\theta_{abs}$) of the motion vectors relative to a common reference plane be computed via the quadrant-corrected arctangent function:

\begin{equation}
\theta_{abs} = \arctan\left(\frac{y_n}{x_n}\right)
\end{equation}

The internal axis angle is extracted from the difference between the absolute angles of the vectors, allowing the dimensionless geometric correction factor ($I_{skew}$) to be isolated as follows:
\begin{equation}
I_{skew} = \tan(\theta_{skew}) = \tan(90^{\circ} - \theta_{int})
\end{equation}

where $x_n$ and $y_n$ are the measured position coordinates along the respective axes. $I_{skew}$ specifies the cross-talk adjustment or parasitic drift per unit displacement required to implement linear software pre-compensation.